% Part: proof-theory
% Chapter: natural-deduction
% Section: sequents

\documentclass[../../../include/open-logic-section]{subfiles}

\begin{document}

\olfileid{pt}{nat}{seq}

\olsection{क्रमवर्तींसह नैसर्गिक निगमन: \Log{N2c} आणि~\Log{N2i}}

~\Log{N1c} किंवा
\Log{N1i} मधील
!!^a{proof}~$\delta$
तिचे अंतिम-!!{formula}~$!A$
हे तिच्या उघड्या गृहीतकांपासून
निष्पन्न होते हे दाखवते.
~$\delta$ मध्ये $!B^x$
हे उघडे गृहीतक असलेल्या
!!{formula}s~$!B$
यांचा संच~$\Gamma$
असेल, तर हे
$\delta \Proves \Gamma
\Sequent !A$ असे
लिहिता येते.
यावरून क्रमवर्तींसाठीही
नैसर्गिक निगमन पद्धत
बांधता येईल असे दिसते.
अशा पद्धतीत प्रत्येक
क्रमवर्ती तिच्या उजवीकडील
!!{formula} कोणत्या
उघड्या गृहीतकांवर अवलंबून
आहे, म्हणजे त्या
क्रमवर्तीत संपणाऱ्या
उप-!!{proof} मध्ये
कोणती गृहीतके उघडी
आहेत, हे सांगते.
यालाच क्रमवर्तींसह,
किंवा “क्रमवर्ती-पद्धतीचे”
नैसर्गिक निगमन म्हणतो.
मिळणाऱ्या पद्धती
\Log{N2c} आणि~\Log{N2i}
अशा आहेत.

~\Log{N1} आणि
\Log{N2} यांतील
संबंध अधिक निकट करण्यासाठी
डावीकडील संच~$\Gamma$
हे खूण लावलेल्या
!!{formula}s~$x:!B$
यांचे धरू. ~\Log{N1}
चे नियम केवळ
आधारविधानांतील
!!{formula}s वर,
म्हणजे निकटच्या
उप-!!{proof}s च्या
अंतिम-!!{formula}s वर
लागतात. त्यामुळे
\Log{N2} चे नियम
क्रमवर्तीच्या उजव्या
भागावरच लागतात.
क्रमवर्तींच्या डावीकडील
खूण लावलेल्या सूत्रांना
म्हणून फक्त संदर्भ
म्हणता येते.

~\Log{N2} सिद्धतांत
गृहीतकांच्या जागी
सर्वात वर अशा रूपाची
स्वयंसिद्धक-क्रमवर्ती असते:
\[x:!A \Sequent !A.\]
नियम~\Log{N1}
च्या नियमांशी तंतोतंत
जुळतात; ते
\olref{tab:N2} मध्ये आहेत.

\subfile{rules-N2}

~\Log{N1c} आणि
\Log{N1i} मध्ये
\Intro{\lif},
\Intro{\lnot},
\Elim{\lor} आणि
\Elim{\lexists} हे
नियम गृहीतके मुक्त करू
शकतात, पण तसे करणे
बंधनकारक नाही.
म्हणून~\Log{N2c}
आणि~\Log{N2i}
मध्येही संबंधित आधारविधानाच्या
संदर्भात अपेक्षित
!!{formula}s $!A$,
$!B$ आणि $!A(a)$
नसली तरी त्या
नियमांचा योग्य उपयोग
मान्य करतो.

\begin{prop}
$\delta$ ही~\Log{N1c}
किंवा \Log{N1i}
मध्ये~$!A$ ची
!!a{proof} असेल,
तर~\Log{N2c}
(\Log{N2i}) मध्ये
$\Gamma \Sequent !A$
ची !!a{proof}~$\delta'$
असते. येथे~$\Gamma$
हा असा सर्व~$x:!B$
यांचा संच आहे की
$!B^x$ हे~$\delta$
चे उघडे गृहीतक असते.
\end{prop}

\begin{proof}
  $n=\pheight{\delta}$
  वर विगमन करू.
  $n=0$ असेल, तर
  $\delta$ ही एकच
  उघडे गृहीतक~$!A^x$.
  मग $\Gamma = \{x:!A\}$
  आणि $\Gamma
  \Sequent !A$ हे
  \Log{N2c} किंवा
  ~\Log{N2i} चे
  स्वयंसिद्धक आहे.

  $n>0$ असेल, तर
  ~$\delta$ मधील अंतिम
  अनुमानानुसार प्रसंग
  पाडू. येथे एकच प्रसंग
  पाहू; उरलेले सरावासाठी.

  अंतिम अनुमान
  $\Intro{\lif}$ असल्याचे
  समजा. मग~$\delta$
  चा शेवट असा होतो:
  \[
  \AxiomC{$\Discharge{!B}{x}$}
  \RightLabel{$\delta_1$}
  \DeduceC{$!C$}
  \DischargeRule{\Intro{\lif}}{x}
  \UnaryInfC{$!B \lif !C$}
  \DisplayProof
  \]
  विगमन गृहीतकानुसार
  ~\Log{N2c} (\Log{N2i})
  मध्ये
  $\Gamma' \Sequent !C$
  ची !!a{proof}~$\delta_1'$
  आहे; $\Gamma'$ म्हणजे
  ~$\delta_1$ मधील
  खूण लावलेली उघडी
  गृहीतके. $!B^x$
  हे $\delta_1$
  चे उघडे गृहीतक असेल,
  तर \Intro{\lif}
  अनुमानात ते मुक्त होते,
  आणि~$\delta$ ची
  उघडी गृहीतके
  $\Gamma =
  \Gamma' \setminus \{x:!B\}$
  होतात. म्हणजे
  ~$\delta_1'$ ही
  $x:!B, \Gamma \Sequent !C$
  ची !!a{proof}
  असते; मग
  $\delta'$ पुढीलप्रमाणे
  घेता येते:
  \[
  \AxiomC{}
  \RightLabel{$\delta_1'$}
  \Deduce$x:!B, \Gamma \fCenter !C$
  \RightLabel{\Intro{\lif}}
  \UnaryInf$\Gamma \fCenter !B \lif !C$
  \DisplayProof
  \]
  $!B^x$ हे~$\delta_1$
  चे उघडे गृहीतक
  नसेल, तर
  \Intro{\lif}
  कोणतेही गृहीतक
  मुक्त करत नाही;
  ~$\delta$ आणि
  $\delta_1$ यांची
  उघडी गृहीतके तीच
  राहतात. ~\Log{N2c}
  आणि~\Log{N2i}
  मध्ये \Intro{\lif}
  च्या आधारविधानाच्या
  संदर्भात $x:!B$
  असणे आवश्यक नसल्याने
  $\delta'$ पुढीलप्रमाणे
  घेता येते:
  \[
  \AxiomC{}
  \RightLabel{$\delta_1'$}
  \Deduce$\Gamma \fCenter !C$
  \RightLabel{\Intro{\lif}}
  \UnaryInf$\Gamma \fCenter !B \lif !C$
  \DisplayProof
  \]

  ~$\delta$ चा शेवट
  इतर नियमांत होणारे
  प्रसंगही असेच आहेत.
\end{proof}

\begin{prop}
$\delta$ ही \Log{N2c}
किंवा \Log{N2i} मध्ये
$\Gamma \Sequent !A$
ची !!a{proof} असेल,
तर~\Log{N1c}
(\Log{N1i}) मध्ये
अंतिम-!!{formula}~$!A$
आणि प्रत्येक
$x:!B \in \Gamma$
साठी उघडे गृहीतक~$!B^x$
असलेली !!a{proof}~$\delta'$
मिळते.
\end{prop}

\begin{proof}
  पुन्हा~$\delta$
  ची !!{height}~$n$
  यावर विगमन करू.
  $n=0$ असेल, तर
  $\delta$ हे
  $x:!A \Sequent !A$
  स्वयंसिद्धक असते.
  !!{proof}~$\delta'$
  म्हणजे गृहीतक~$!A^x$.

  $n>0$ असेल,
  तर~$\delta$ मधील
  अंतिम अनुमानानुसार
  प्रसंग पाडू.
  उदाहरणार्थ, ते
  ~$\Intro{\lif}$
  असेल, तर~$\delta$
  चा शेवट असा:
  \[
  \bottomAlignProof
  \AxiomC{}
  \RightLabel{$\delta_1$}
  \Deduce$x:!B, \Gamma \fCenter !C$
  \RightLabel{\Intro{\lif}}
  \UnaryInf$\Gamma \fCenter !B \lif !C$
  \DisplayProof
  \quad\text{ किंवा }\quad
  \bottomAlignProof
  \AxiomC{}
  \RightLabel{$\delta_1$}
  \Deduce$\Gamma \fCenter !C$
  \RightLabel{\Intro{\lif}}
  \UnaryInf$\Gamma \fCenter !B \lif !C$
  \DisplayProof
  \]
  विगमन गृहीतकानुसार
  ~\Log{N1c} (\Log{N1i})
  मध्ये~$!C$ ची
  !!a{proof} मिळते.
  पहिल्या प्रसंगी
  $!B^x$ हे~$\delta_1'$
  चे उघडे गृहीतक असते;
  दुसऱ्या प्रसंगी नसते.
  मग !!{proof}~$\delta'$
  अनुक्रमे पुढीलप्रमाणे
  घेता येते:
  \[
  \bottomAlignProof
  \AxiomC{$\Discharge{!B}{x}$}
  \RightLabel{$\delta_1'$}
  \DeduceC{$!C$}
  \DischargeRule{\Intro{\lif}}{x}
  \UnaryInfC{$!B \lif !C$}
  \DisplayProof
  \quad\text{ किंवा }\quad
  \bottomAlignProof
  \AxiomC{}
  \RightLabel{$\delta_1'$}
  \DeduceC{$!C$}
  \RightLabel{\Intro{\lif}}
  \UnaryInfC{$!B \lif !C$}
  \DisplayProof
  \]
\end{proof}

\end{document}
